\chapter{Correspondencia con las lecturas}
Las 32 lecturas del primer curso se integran en \textbf{seis partes y dieciséis capítulos temáticos}. Se incluyen los \textbf{206 ejercicios originales numerados}, con sus incisos y soluciones. La lectura 9 contiene un ejercicio bajo un encabezado singular; la lectura 25 no tiene una lista independiente. Las verificaciones propuestas como ejercicio dentro de la teoría también se desarrollan en el texto.

Cada ejercicio identifica su lectura y número originales entre paréntesis. Su número principal sigue la numeración del capítulo del libro. La tabla permite localizar el comienzo del desarrollo de cada lectura; cuenta ejercicios numerados, no incisos ni ejemplos añadidos. Cuando un enunciado requiere corrección, el texto explica el problema y demuestra una versión válida.

Los resultados geométricos de la lectura 22 se estudian en el capítulo~\ref{cap:nuevo08}; su criterio para operadores de rango finito y los ejercicios~\ref{ej:lectura-22-1}, \ref{ej:lectura-22-2}, \ref{ej:lectura-22-3} y~\ref{ej:lectura-22-6} se sitúan en el capítulo~\ref{cap:nuevo09}, después de introducir la norma de operador. Las herramientas sobre convergencia uniforme usadas por Tietze se desarrollan en el capítulo~\ref{cap:nuevo11} y se aplican en el capítulo~\ref{cap:nuevo13}.

El criterio que caracteriza los funcionales continuos por su núcleo, procedente de la lectura 21, se demuestra en la proposición~\ref{ch21:kernel-cerrado} del capítulo~\ref{cap:nuevo09}, antes de su aplicación al rango finito. El desarrollo del dual y los ejercicios de esa lectura permanecen en el capítulo~\ref{cap:nuevo10}.

\footnotesize
\begin{longtable}{@{}r>{\raggedright\arraybackslash}p{.50\textwidth}rrr@{}}
\toprule
Lectura & Tema de origen & Cap. & Ej. & Pág.\\\midrule
\endfirsthead
\toprule Lectura & Tema de origen & Cap. & Ej. & Pág.\\\midrule
\endhead
\bottomrule\endfoot
1 & Distancias y normas & \ref{cap:nuevo01} & 7 & \pageref{lectura:01}\\
2 & Bolas, vecindades, interior y clausura & \ref{cap:nuevo01} & 7 & \pageref{lectura:02}\\
3 & Abiertos, cerrados y sucesiones & \ref{cap:nuevo01} & 9 & \pageref{lectura:03}\\
4 & La topología de un subespacio & \ref{cap:nuevo01} & 6 & \pageref{lectura:04}\\
5 & Continuidad y equivalencia topológica & \ref{cap:nuevo02} & 9 & \pageref{lectura:05}\\
6 & Densidad y numerabilidad & \ref{cap:nuevo02} & 8 & \pageref{lectura:06}\\
7 & Continuidad uniforme y control de las distancias & \ref{cap:nuevo02} & 7 & \pageref{lectura:07}\\
8 & Espacios completos y espacios de Banach & \ref{cap:nuevo03} & 7 & \pageref{lectura:08}\\
9 & La completación de un espacio métrico & \ref{cap:nuevo03} & 1 & \pageref{lectura:09}\\
10 & Categoría de Baire y consecuencias de la completitud & \ref{cap:nuevo04} & 10 & \pageref{lectura:10}\\
11 & El principio de contracción y sus aplicaciones & \ref{cap:nuevo04} & 6 & \pageref{lectura:11}\\
12 & Separación de cerrados y extensión de Tietze & \ref{cap:nuevo13} & 3 & \pageref{lectura:12}\\
13 & Compacidad: del control local al control global & \ref{cap:nuevo05} & 8 & \pageref{lectura:13}\\
14 & Compacidad y sucesiones & \ref{cap:nuevo05} & 6 & \pageref{lectura:14}\\
15 & Acotación total y compacidad local & \ref{cap:nuevo05} & 6 & \pageref{lectura:15}\\
16 & Compactificaciones y espacios universales & \ref{cap:nuevo07} & 5 & \pageref{lectura:16}\\
17 & Conexidad y continuidad & \ref{cap:nuevo06} & 7 & \pageref{lectura:17}\\
18 & Componentes conexas, caminos y arcos & \ref{cap:nuevo06} & 7 & \pageref{lectura:18}\\
19 & Geometría y series en espacios normados & \ref{cap:nuevo08} & 7 & \pageref{lectura:19}\\
20 & Operadores acotados y acotación uniforme & \ref{cap:nuevo09} & 7 & \pageref{lectura:20}\\
21 & El dual, los isomorfismos y la extensión de funcionales & \ref{cap:nuevo10} & 8 & \pageref{lectura:21}\\
22 & Qué distingue a la dimensión finita & \ref{cap:nuevo08}, \ref{cap:nuevo09} & 6 & \pageref{lectura:22}, \pageref{ch22:rango-finito-kernel}\\
23 & Operadores compactos & \ref{cap:nuevo09} & 4 & \pageref{lectura:23}\\
24 & Convexidad, separación y puntos extremos & \ref{cap:nuevo10} & 7 & \pageref{lectura:24}\\
25 & La métrica uniforme y los espacios de funciones & \ref{cap:nuevo11} & 0 & \pageref{lectura:25}\\
26 & Convergencia uniforme: límites, series, integrales y derivadas & \ref{cap:nuevo11} & 8 & \pageref{lectura:26}\\
27 & Equicontinuidad y el teorema de Arzelà--Ascoli & \ref{cap:nuevo12} & 5 & \pageref{lectura:27}\\
28 & Álgebras de funciones y aproximación de Stone--Weierstrass & \ref{cap:nuevo13} & 4 & \pageref{lectura:28}\\
29 & Producto interno y geometría de Hilbert & \ref{cap:nuevo14} & 8 & \pageref{lectura:29}\\
30 & Proyecciones ortogonales y mejor aproximación & \ref{cap:nuevo14} & 6 & \pageref{lectura:30}\\
31 & Bases ortonormales y desarrollos de Fourier & \ref{cap:nuevo15} & 7 & \pageref{lectura:31}\\
32 & Representación de Riesz y operadores adjuntos & \ref{cap:nuevo16} & 10 & \pageref{lectura:32}\\
\end{longtable}
\normalsize
